\section*{Chapter \ref{ch:b2:diffraction} --- Fraunhofer Diffraction and Spatial Filtering}

\begin{solution}{exo:b2:diffraction:1}
$2\lambda L/b = \qty{25}{mm}$; secondary maxima at $1.43\lambda L/b = \qty{18}{mm}$ and
$2.46\lambda L/b = \qty{31}{mm}$; \omterm{rem:b2:diffraction:fresnelnumber}{Fresnel number} $b^2/\lambda L = 8 \times 10^{-3}$: far
enough. With the lens, $2\lambda f/b = \qty{6.3}{mm}$.
\end{solution}

\begin{solution}{exo:b2:diffraction:2}
$1.22\lambda/D$: eye \qty{2.2e-4}{rad} (\qty{46}{''}); \qty{10}{cm}: \qty{1.4}{''}; Hubble
\qty{0.058}{''}; dish $1.22 \times 0.06/64 = \qty{1.1e-3}{rad} = \qty{3.9}{'}$; camera:
$\qty{1.1e-4}{rad} \times \qty{50}{mm} = \qty{5.6}{\micro m}$ radius, an \qty{11}{\micro m} disc ---
three pixels: \omterm{def:b2:diffraction:huygens}{diffraction} already limits a sensor at $f/8$.
\end{solution}

\begin{solution}{exo:b2:diffraction:3}
$\lambda f/a = \qty{2}{mm}$; envelope $2\lambda f/b = \qty{20}{mm}$; ten fringes; orders
$p = 5, 10, \dots$ ($a/b = 5$) missing.
\end{solution}

\begin{solution}{exo:b2:diffraction:4}
$\theta = 1.22\lambda/D = \qty{3.9e-4}{rad}$: a spot \qty{290}{km} across; with \qty{3.5}{m},
\qty{170}{m}. The reflector's own \omterm{def:b2:diffraction:huygens}{diffraction}, $\lambda/\qty{1}{m}$, spreads the
return over \qty{600}{m}: of $10^{17}$ photons sent, a handful come back.
\end{solution}

\begin{solution}{exo:b2:diffraction:5}
(a) $\int_{-b/2}^{b/2}\eu^{\iu kx\sin\theta}\dd x = b\,\operatorname{sinc}(kb\sin\theta/2)$. (b)
$\dd(\sin x/x)/\dd x = 0$ gives $\tan x = x$; $x = 4.49 = 1.43\pi$, $\operatorname{sinc}^2 =
0.047$. (c) $0.903$. (d) The phase becomes $kx(\sin\theta - \sin\theta_0)$: the
same pattern about $\theta_0$.
\end{solution}

\begin{solution}{exo:b2:diffraction:6}
$0.61\lambda/\mathrm{NA}$: \qty{220}{nm}; \qty{170}{nm}; \qty{110}{nm}. Beyond some $500\times$
the eyepiece only enlarges the blur (empty magnification). Electrons
at \qty{4}{pm} with $\mathrm{NA} \approx 0.01$: \qty{0.2}{nm} --- atoms.
\end{solution}

\begin{solution}{exo:b2:diffraction:7}
(a) The integral factorizes in $x$ and $y$. (b) $2\lambda f/b = \qty{5}{mm}$. (c) A
disc's chords are shorter than its diameter away from the centre: its
effective width is smaller, its pattern wider. (d) Four \omterm{def:b2:diffraction:huygens}{diffraction}
spikes in a cross.
\end{solution}

\begin{solution}{exo:b2:diffraction:8}
(a) $\sin\theta = \pm0.275$ (\qty{16}{\degree}), $\pm0.55$ (\qty{33}{\degree}). (b) $\sin\alpha \ge 0.275$.
(c) Orders $0$, $\pm1$ only: a sinusoidal image of the right period, not
the sharp bars. (d) $\sin\alpha \ge \lambda/2d = 0.14$: twice the resolution.
\end{solution}

\begin{solution}{exo:b2:diffraction:9}
(a) $2.44\lambda f/D = \qty{15}{\micro m}$; pinhole \qty{23}{\micro m}. (b) $\lambda/\qty{50}{\micro m} =
\qty{1.3e-2}{rad}$: \qty{0.25}{mm} off axis, blocked. (c) $1\%$; the pinhole keeps only
the low spatial frequencies. (d) Offset by a few micrometres and the
pinhole clips the spot itself.
\end{solution}

\begin{solution}{exo:b2:diffraction:10}
(a) Amplitudes add to the unobstructed wave, which is zero off axis:
$|\underline s_{\text{slit}}| = |\underline s_{\text{strip}}|$ there. (b) Fringes $\lambda L/b =
\qty{16}{mm}$ apart; measure the spacing, get $b$. (c) The strip's shadow sits
in the bright undiffracted beam. (d) Droplets of diameter $d$ diffract
like holes: rings at $1.22\lambda/d$, red outside; the radius gives $d$.
\end{solution}

\begin{solution}{exo:b2:diffraction:11}
(a) $|1 + \iu\varphi|^2 = 1 + \varphi^2$: uniform to first order. (b) Without the
axial light, $|\iu\varphi|^2 = \varphi^2$: second order, faint. (c) $|\iu + \iu\varphi|^2 =
1 + 2\varphi + \dots$: linear. (d) The axial spot outshines the diffracted
light; attenuating it evens the two amplitudes and raises the contrast.
\end{solution}

\begin{solution}{exo:b2:diffraction:12}
(a) Half-width $\lambda/W$; divided by $p/a\cos\theta$: $\Delta\lambda = \lambda a\cos\theta/pW \approx
\lambda/pN$. (b) The collimated beam has the diameter $Df_{\text{coll}}/f_{\text{tel}}$;
a smaller grating would waste light and resolving power. (c) A
resolution is always the angle over which the phase across the
aperture changes by one wavelength, $\lambda/(\text{size})$. (d) Its image
width, which dominates when wider than the \omterm{def:b2:diffraction:huygens}{diffraction} width of the
grating --- the usual case.
\end{solution}

\begin{solution}{pb:b2:diffraction:1}
\textbf{1.} $1.22\lambda/D = \qty{2.2e-4}{rad}$; $\times\qty{17}{mm} = \qty{3.8}{\micro m}$: two cone
spacings --- the retina samples at the optics' limit.

\textbf{2.} \qty{56}{\micro m} at \qty{25}{cm}, \qty{2.2}{mm} at \qty{10}{m}; \qty{1}{mm} at
\qty{50}{cm} is ten times the limit: easily.

\textbf{3.} \qty{9.6e-5}{rad}; the lens's aberrations at full aperture and the
coarser, pooled rods undo the gain.

\textbf{4.} $1.5/2.2 \times 10^{-4} = \qty{7}{km}$.

\textbf{5.} \omterm{def:b2:diffraction:huygens}{Diffraction} worsens to \qty{6.7e-4}{rad}, but the blur circle of
a defocused eye shrinks with the aperture: depth of field wins.

\textbf{6.} \qty{6.7e-7}{rad} (\qty{0.14}{''}); $\times\qty{10}{m} = \qty{6.7}{\micro m}$ radius: a
disc of $1.3$ pixels.

\textbf{7.} \qty{1}{''} $= \qty{4.9e-6}{rad}$, seven times the limit; $D = 1.22\lambda/4.9
\times 10^{-6} = \qty{14}{cm}$ already reaches the seeing.

\textbf{8.} \qty{2.7e-6}{rad} (\qty{0.55}{''}); the atmosphere's phase errors are
a smaller fraction of a longer wavelength: fewer, slower corrections.

\textbf{9.} $0.15 < 0.55$: no; at \qty{550}{nm} from space, \qty{0.14}{''}: just.

\textbf{10.} $6.7 \times 10^{-7} \times 3.8 \times 10^8 = \qty{250}{m}$; the lander, no.

\textbf{11.} A cross of four \omterm{def:b2:diffraction:huygens}{diffraction} spikes.

\textbf{12.} $(1000/7)^2 = 2 \times 10^4$: fainter objects, not sharper ones.

\textbf{13.} \qty{370}{nm} and \qty{240}{nm}: $2700$ and $4200$ line pairs per
millimetre.

\textbf{14.} The bacterium, yes; its \qty{0.2}{\micro m} structure at the limit
with oil, not dry; the virus, no.

\textbf{15.} $d \ge \lambda/\mathrm{NA} = \qty{390}{nm}$ under normal illumination, \qty{200}{nm}
with oblique light --- the $0.61\lambda/\mathrm{NA}$ of two points lies between.

\textbf{16.} The oil's index raises $n\sin\alpha$ and suppresses the total
reflection that would trap the steep rays in the coverslip.

\textbf{17.} \qty{110}{nm}; electrons: $0.61 \times 4\,\text{pm}/0.01 = \qty{0.24}{nm}$.

\textbf{18.} The plate retards the undiffracted light by a quarter wave
so that it interferes with the light diffracted by the cell's
phase structure; the phase pattern becomes an intensity pattern.

\textbf{19.} \qty{4.3e-4}{rad}; $1.5 + 2 \times 43 = \qty{88}{mm}$.

\textbf{20.} $2.44\lambda f_1/D = \qty{8.7}{\micro m}$; pinhole \qty{13}{\micro m}.

\textbf{21.} $f_2 = 20f_1 = \qty{200}{mm}$; \qty{2.2e-5}{rad}; $30 + 44 = \qty{74}{mm}$ after a
kilometre.

\textbf{22.} $\lambda/\qty{100}{\micro m} = \qty{5.3e-3}{rad}$: \qty{53}{\micro m} off axis at the
pinhole --- blocked.

\textbf{23.} $1\%$, absorbed by the pinhole plate.

\textbf{24.} A beam of diameter $D$ doubles over $\sim D^2/\lambda$: \qty{4}{m} for
\qty{1.5}{mm}, \qty{1.7}{km} for \qty{30}{mm}.

\textbf{25.} Eye \qty{3}{mm}, telescope \qty{1}{m}, objective (NA $1.4$), beam
\qty{30}{mm}: in each, $\lambda/D$ sets the smallest angle, the finest detail, the
slowest spread.
\end{solution}
